\section{Retrieval-Augmented Generation for Software Tasks}
\label{sec:sota-rag}

Retrieval is the default answer to the repository-context problem, and for good
reason: it needs no language-specific tooling, scales to any file the embedder can
read, and requires nothing of the repository beyond that it be
readable~\cite{lewis2020rag,karpukhin2020dpr}. RepoCoder applies it to
repository-level completion and, notably, does not apply it once: it alternates
retrieval and generation so that partially generated code can re-query the
repository, on the finding that a single query formed from the immediate context
retrieves poorly~\cite{zhang2023repocoder}. GraphRAG responds to the same weakness
from a different direction, building an entity graph over the corpus because some
questions are about how documents connect rather than about any one
document~\cite{edge2024graphrag}.

Both are, read together, admissions about the limits of similarity as a relevance
criterion, and both point at the comparison this thesis makes. If similarity
retrieval were adequate for repository context, neither iterative re-querying nor
graph construction would be necessary. If it were useless, it would not be the
default. The interesting question is what each kind of selection is good \emph{at},
and that question is only answerable by an experiment in which the two are given
the same task, model, evidence pack and evaluation under a matched prompt template.
In this thesis, the rendered context and the instruction sentence that names it
vary by arm, a limitation reported explicitly in Chapter~\ref{chap:results}.

That framing also predicts where the two should differ, which is what makes the
comparison falsifiable rather than exploratory. Similarity retrieval returns code
that resembles the change, which should help a reviewer say something well-informed
about the change's own subject matter, including conventions, comparable
implementations, and local idiom. Structural traversal returns code connected to the change, which
should help specifically with the questions about consequence: what else depends on
this, what tests exercise it, what breaks. A study that reports only a single
aggregate quality score cannot see this distinction even if it is present, and one
of the design decisions in Chapter~\ref{chap:edas} is to report the targeted subset
of criteria separately for exactly that reason.
